\documentclass[12pt]{article}
\usepackage[utf8]{inputenc}
\usepackage[T1]{fontenc}
\usepackage{hyperref}
\title{Karore Language Resources for K-12 Curriculum Development}
\author{Compiled by Claude Sonnet 5.5 (claude-sonnet-5-5) for Andr\'as Kornai}
\date{}
\begin{document}
\maketitle

\section{Introduction}
Karore (ISO 639-3 xkx) is an Austronesian language of West New Britain Province, Papua New Guinea. It belongs to the Pasismanua subgroup of the Arawe-Pasismanua branch of the Ngero-Vitiaz family. It is spoken by roughly 550 people (an Ethnologue figure from 2003) between the Andru and Johanna Rivers. The community is remote. Joshua Project reports that Karore remains the primary language. Endangerment ratings vary by source (UNESCO: Vulnerable). I found no school materials in the search results.

\section{Primary Educational Resources}
None found. I found no primer, textbook or school programme for Karore.

\section{Linguistic and Reference Materials}
\begin{itemize}
\item Wikipedia, ``Karore language'': an overview of classification and location. \url{https://en.wikipedia.org/wiki/Karore_language}
\item Endangered Languages Project entry for Karore. \url{https://endangeredlanguages.com/lang/2224}
\end{itemize}

\section{Religious and Community Resources}
\begin{itemize}
\item Joshua Project profile of the Karore people. It reports that there is no Karore Bible, that translation is needed, and that speakers can partly use the related Kaulong scriptures. \url{https://joshuaproject.net/people_groups/12551}
\end{itemize}

\section{Digital and Audio Resources}
No audio, video or corpus resources were found in the search results.

\section{Curriculum Development Context}
There is almost no documentation beyond basic reference listings. A native-language curriculum would require new documentation work, such as an orthography and a word list, together with the community. Papua New Guinea's elementary vernacular-education approach is a possible model, but I did not verify any Karore involvement in it.
\end{document}